\section{Approximation from smoothed marginal states}

The marginal states at different cuts may be smoothed independently.
Throughout this section the smoothing radius is measured in the full trace norm.

\begin{theorem}[Truncation after independent smoothing]
    \label{thm:smoothed_marginal_truncation}
    \lean{TNLean.MPS.AreaLaw.norm_sequentialTruncation_leadingEigenprojection_sub_sq_le_of_smoothing}
    \leanok
    Let $\psi$ be a unit vector on a chain of $N$ sites, let $D\in\mathbb N$, and
    let $\rho_k$ be its marginal on the first $k$ sites. Suppose that for each
    internal cut $1\le k<N$ there is a density matrix $\sigma_k$ such that
    \[
        \|\rho_k-\sigma_k\|_1\le\eta_k,
        \qquad \tau_D(\sigma_k)\le T_k,
    \]
    where $\tau_D$ is the sum of the eigenvalues after the first $D$ in
    decreasing order. Let $\Phi_D$ be the sequential truncation of $\psi$
    using the leading $D$ eigenvectors of each $\rho_k$. Then
    \[
        \|\Phi_D-\psi\|_2^2
        \le \sum_{k=1}^{N-1}\left(T_k+\frac{\eta_k}{2}\right).
    \]
    This gives the finite-chain estimate for the smoothing argument
    of~\cite{SchuchWolfVerstraeteCirac2007}.
\end{theorem}
\begin{proof}
    \leanok
    \uses{thm:chain_sequential_truncation_sq}
    Since both marginal states have trace one, the variational characterization
    of the leading spectral sum gives
    \[
        \tau_D(\rho_k)\le\tau_D(\sigma_k)
        +\frac12\|\rho_k-\sigma_k\|_1.
    \]
    Sum this inequality over the internal cuts and apply the squared-error
    truncation bound. No common smoothing state is required.
\end{proof}
